\subsection{Motivierendes Beispiel}
\begin{frame}[fragile]{Prädikatenlogik among us}
    \begin{itemize}
        \item<1-> Ihr seid auf einem Raumschiff
        \item<2-> Mit eurer Crew 
        $$
        U = \{\sus{red}, \sus{green}, \sus{blue}, \sus{white}, \sus{yellow}, \sus{magenta}\}
        $$
        \item<3-> Aber manche von euch sind Imposter
        
    \end{itemize}
    \only<2>{\hspace{3cm}\includegraphics[width=4.5cm]{../figures/amongus_crew.png}} 
    \only<3>{\hspace{3cm}\includegraphics[width=4.5cm]{../figures/amongus_imposter.jpg}} 
\end{frame}

\begin{frame}[fragile]{Prädikatenlogik among us}
    \begin{itemize}
        \item<1-> Wir brauchen eine Funktion um Imposter zu erkennen
        $$C: U \to \{wahr, falsch\}$$
        \item<2-> Und eine Funktion um zu beschreiben ob ein Crewmitglied noch lebt
        $$A: U \to \{wahr, falsch\}$$
        \item<3-> Und eine Funktion um zu beschreiben was ein Imposter macht wenn er ein Crewmitglied findet
        $$K: U\times U \to \{wahr, falsch\}$$
    \end{itemize}
\end{frame}


\begin{frame}[fragile]{Eine kurze Geschichte}
$$
\overline{C(\sus{red})}\wedge C(\sus{white})\wedge K(\sus{red},\sus{white})\implies \overline{A(\sus{white})}
$$
\end{frame}

\begin{frame}[fragile]{Aber wer ist den Imposter?}
\begin{itemize}
    \item<1-> \alert{Problem:} Wer ist den Imposter?
    \item<2-> Wir wissen nicht welche Crewmitglieder wir für was einsetzen können 
    \item<3-> Also müssen wir allgemeiner werden:
    $$
    \exists x,z \in U: \overline{C(x)}\wedge C(y)\wedge K(x,y)\implies \overline{A(y)}
    $$
\end{itemize}
\begin{center}
\only<4>{Das ist eine prädikatenlogische Formel}
\end{center}
\end{frame}

\subsection{Definitionen}

\begin{frame}{Prädikatenlogik erster Ordnung}
\begin{itemize}
    \item<1-> \alert{Variablen} sind wie ihr sie aus der Aussagenlogik bereits kennt.
    \item<2-> $U$ ist unser \alert{Universum}. Es enthält alle Elemente die wir in unsere Variablen einsetzen dürfen 
    \item<3-> \alert{Prädikatensymbole} sind Mappings von Elementen aus $U$ auf $wahr$ oder $falsch$.
    Sie werden normal groß geschrieben.
    \item<4-> \alert{Funktionssymbole} sind Mappings von Elementen aus $U$ auf andere Elemente aus $U$.
    Sie werden normal klein geschrieben.
    \item<5> Eine \alert{Interpretation} $I$ ist ein Lösungsvorschlag für die prädikatenlogische Formel.
    Es ist ein Mapping von Variablen auf Elemente des Universums, der Prädikatensymbole auf Prädikate und der Funktionsformeln auf Funktionen.
\end{itemize}
\end{frame}

\begin{frame}{Unser Beispiel}
    \begin{itemize}
        \item $U = \{\sus{red}, \sus{green}, \sus{blue}, \sus{white}, \sus{yellow}, \sus{magenta}\}$
        \item Prädikate:
        \begin{itemize}
            \item $P_{sus}(x,y)$: $x$ verdächtigt $y$
            \item $P_{kill}(x,y)$: $x$ tötet $y$
            \item $P_{seen}(x,y)$: $x$ und $y$ wurden zusammen gesehen
            \item $P_{alive}(x)$: $x$ lebt (noch)
            \item $P_{crew}(x)$: $x$ ist ein echtes Crewmitglied
        \end{itemize}
    \end{itemize}
\end{frame}

{\setbeamercolor{palette primary}{bg=ExColor}
\begin{frame}[fragile]{Denkpause}
    \footnotesize
        \begin{alertblock}{Aufgaben}
            Findet Interpreatationen für die folgenden Prädikatenlogischen Formeln in unserem Beispiel:
        \end{alertblock}
        \metroset{block=fill}
        \begin{block}{Normal}
            \begin{itemize}
                \item $(P_1(x,y) \implies P_2(x)) \implies P_3(y)$
                \item $(P_1(x,y) \implies P_2(y)) \implies P_3(x)$
            \end{itemize}
        \end{block}
        \begin{block}{Etwas Schwerer}
            \begin{itemize}
                \item $(P_1(x,y) \wedge P_2(y)) \implies P_3(y)$
            \end{itemize}
        \end{block}
\end{frame}
}

\begin{frame}{Quantoren und Funktionen}
    \begin{itemize}
        \item TODO : Quantoren und Funktionen erklären, dann auf Mathematische Beispiele wechseln (funktionen +, - , * , / und U = N...)
    \end{itemize}
\end{frame}

\begin{frame}{Ein weiteres Beispiel}
    \begin{itemize}
        \item $U = \mathbb Z$
        \item Prädikate:
        \begin{itemize}
            \item $P_{even}(x)$: $x$ is even
            \item $P_{prim}(x)$: $x$ is prim
            \item $P_{+}(x)$: $x$ is positive
            \item $P_>(x,y): x>y$ 
            \item $P_=(x,y): x=y$ 
        \end{itemize}
        \item Funktionen: $+, -, \cdot$ 
    \end{itemize}
\end{frame}

{\setbeamercolor{palette primary}{bg=ExColor}
\begin{frame}[fragile]{Denkpause}
    \footnotesize
        \begin{alertblock}{Aufgaben}
            Findet Interpreatationen für die folgenden Prädikatenlogischen Formeln in unserem Beispiel:
        \end{alertblock}
        \metroset{block=fill}
        \begin{block}{Normal}
            \begin{itemize}
                \item $\exists x,y : P_1(x) \wedge P_1(y) \implies \overline{P_1(f_1(x,y))}$
                \item $\forall x\exists y: P_1(f_1(x,y),x)$
                \item $\forall x\exists y: P_1(f_1(x,y),y)$
            \end{itemize}
        \end{block}
        \begin{block}{Etwas Schwerer}
            \begin{itemize}
                \item TODO
            \end{itemize}
        \end{block}
\end{frame}
}